\documentclass[12pt]{IEEEtran}
\usepackage{cite}
\usepackage{amsmath,amssymb,amsfonts}
\usepackage{algorithmic}
\usepackage{graphicx}
\usepackage{textcomp}
\def\BibTeX{{\rm B\kern-.05em{\sc i\kern-.025em b}\kern-.08em
    T\kern-.1667em\lower.7ex\hbox{E}\kern-.125emX}}
\begin{document}
\title{Human Activity Recognition via Unsupervised Clustering Methods}
\author{CS 472/572, Spring 2026 \\ Bailey Hall and Josselin Retiguin}

\maketitle

\section{Introduction}
\label{sec:intro}

The ever-increasing adoption of wearable devices and smartphones has allowed for the integration of high-frequency inertial sensors to become a part of many people's daily lives. This has led Human Activity Recognition (HAR) to transform from a niche research area into a primary focus in mobile health, elderly care, and personalized fitness. The sheer volume of data generated by these sensors has steered research toward robust unsupervised clustering approaches that can interpret human movement without human intervention. However, processed sensor data is inherently high-dimensional and often contains complex, non-linear relationships between data features that linear dimensionality reduction techniques used in traditional HAR clustering research fail to capture. 

Previous HAR literature has explored a wide range of clustering algorithms, yet few papers have investigated which specific models and approaches are most effective for more focused, single-dataset applications. \textbf{The objective of this project is to provide a comparative analysis of unsupervised clustering performance on the UCI Human Activity Recognition Using Smartphones Dataset.} Rather than compare a limited number of algorithms across multiple datasets, this study focuses on five distinct clustering models: K-Means, Hierarchical Agglomerative, Spectral, DBSCAN, and Gaussian Mixture Models, each applied to a single standardized dataset.

In addition to evaluating traditional dimensionality reduction techniques like Principal Component Analysis (PCA), we introduce an extension into neural network-based feature extraction using a symmetric Autoencoder (AE) and a Convolutional Variational Autoencoder (CVAE). By compressing the original 561-dimensional sensor data into a 50-dimensional latent space, this project investigates how non-linear and spatial-temporal feature learning can impact clustering performance.

The results detailed in this paper demonstrate that while traditional unsupervised approaches can achieve a respectable baseline, the use of CVAE-based feature extraction can significantly improve clustering outcomes, achieving an Adjusted Rand Index (ARI) score approximately 54\% higher than the linear baseline. This project highlights the potential for hybrid deep learning + unsupervised learning frameworks to move research toward more accurate and efficient human activity recognition.

\section{Literature Review}
\label{sec:lit}
Advances in machine learning have shifted Human Activity Recognition (HAR) from using labor-intensive supervised learning models for the classification of human activities to unsupervised clustering frameworks. As the amount of data from wearable sensors increases, the traditional method of manually annotating data has become less feasible. A review of clustering techniques in HAR [1] identifies four phases in the unsupervised learning process consisting of feature extraction, algorithm selection, cluster validation, and interpretation. While state-of-the-art research in HAR often utilizes more complex neural networks for classification, classic techniques such as K-Nearest Neighbors and sub-clustering continue to prove both useful and efficient, with some implementations achieving over 98\% on certain datasets [1].

An important aspect of successful unsupervised learning is choosing which clustering method best fits the dataset and the information known about the dataset beforehand. For example, in an HAR research study that utilized smartphone sensors [2], the researchers found that Gaussian Mixture Models (GMM) can achieve 100\% accuracy in classifying physical activities when the number of target clusters is known and given to the model ahead of time. However, when no pre-defined number of clusters is passed to the model, algorithms such as Density-Based Spatial Clustering of Applications with Noise (DBSCAN) [2][3][4] and average-linkage hierarchical agglomerative clustering (HIER) [2] demonstrate higher clustering performance. These algorithms often utilize metrics such as the Calinski-Harabasz (CH) index [2] to automatically calculate the optimal number of clusters, which removes the need for researchers to know the target number of clusters beforehand.

As data gathered from wearable sensors becomes increasingly high-dimensional, the challenges of both accuracy and computational efficiency become more difficult to address. Research on feature selection in high-dimensional activity data has shown that identifying and removing redundant data features with Correlation Feature Selection (CFS) can significantly improve cluster separability and reduce processing time up to 71\% [3]. Similarly, using dimensionality reduction techniques like Principal Component Analysis (PCA) and Uniform Manifold Approximation and Projection (UMAP)  allows for the capture of complex human movement patterns while increasing the efficiency of models like Spectral Clustering [3][4][5]. Further, one study [5] found that rather than considering each feature corresponding to individual tri-axial data points, focusing on the Acceleration Vector Magnitude (AVM) yields more accurate clustering results. 

More recent research in HAR involves hybrid approaches that combine deep learning and unsupervised clustering. One paper [6] introduces one such method, an adapted Semantic Clustering by Adopting Nearest Neighbors (SCAN) framework. The researchers found that SCAN, which combines the self-supervised pretext task SimCLR with nearest-neighbor clustering, outperformed state-of-the-art deep sensory clustering across multiple benchmarks [6]. These hybrid approaches are particularly useful for clustering imbalanced datasets and sporadic activities, which frequently occur when gathering data from naturalistic human environments [6].

In addition to general activity recognition, HAR unsupervised clustering methods have more specialized safety applications. An example of this is fall detection in elderly patients, where GMM and Fuzzy C-Means (FCM) have been shown to distinguish between daily-life activities and simulated falls with up to 89\% accuracy [5]. Another application is in driving safety and vehicle accident prevention. One study [7] applied K-means clustering to tri-axial acceleration and rotation data from smartphone sensors in order to predict driver behavior with an achieved accuracy of approximately 86\%. 

\begin{figure}[!t]
\centerline{\includegraphics[width=\columnwidth]{lit review summary.png}}
\caption{Summary Table of the Clustering Models used in the Literature Reviewed. The top-performing models found in each paper are listed in the "Overall Best" row of the table.}
\label{fig1}
\end{figure}

This review of recent unsupervised clustering applications in HAR research demonstrates that while the field is moving toward deep learning and neural network approaches [6], a combination of efficient feature selection techniques [3] and well-established clustering algorithms [2][5][7] remains a practical and useful approach for real-world activity recognition. \textbf{Fig. 1} provides a summary of models used for clustering in the literature reviewed.

\section{Methods}
The majority of the literature reviewed utilized data gathered from a wearable device and/or smartphone, specifically data sensed from a device's built-in accelerometer and gyroscope. In addition, multiple studies used and compared different clustering models on multiple datasets. 

In order to continue the work of the literature, the goal of this project is to compare a range of common unsupervised clustering algorithms on data from wearable sensors. \textbf{However, rather than compare 1-3 algorithms to each other on multiple datasets, we aim to compare the performance of 5 different clustering algorithms on a single dataset.} This is a more focused approach that can demonstrate the best clustering approach for datasets of the same type and format as the dataset we are working with.

\subsection{Dataset and Preprocessing}
The dataset utilized in this study is the \textbf{UCI Human Activity Recognition Using Smartphones Dataset} [8]. The data in this dataset were collected from 30 subjects wearing a Samsung Galaxy S II smartphone on their waist. The subjects then performed six distinct activities: laying, sitting, standing, walking, walking downstairs, and walking upstairs. Each record in the dataset consists of 561 features extracted from the raw tri-axial acceleration and tri-axial velocity data gathered by the smartphone's embedded accelerometer and gyroscope.

The UCI HAR dataset was chosen in part due to the fact that accelerometers and gyroscopes are among the most common sensors an average person carries with them during daily life. In addition to being built into the majority of smartwatches and smart rings, as the name of the dataset suggests, both sensors are also in all modern smartphones. This means that specialized equipment is not needed in order to gather more data of the same type and format as this dataset, and that any HAR discoveries made using this data have the potential to improve the health of anyone who owns a basic mobile smart device.

The data is provided in a .txt format, with a 70-30 train-test split. Additionally, for both the training and testing data, the subject data, labels, and features are all in separate files. To better align with the unsupervised characteristic of this project and for ease of manipulation, all of the separate data files were combined into a single dataframe, maximizing the amount of data that would be used to learn the clusters. Next, after extracting the labels from the dataframe, each feature in the dataset was standardized using Scikit-learn's \textbf{StandardScaler} in order to ensure that high-magnitude features did not skew any distance-based algorithms. The result of preprocessing the data is a NumPy array with 10,299 records, each with 561 features, where each feature has a distribution with a mean of 0 and a standard deviation of 1.

\subsection{Dimensionality Reduction}
As was highlighted in some of the literature reviewed [3][5], high-dimensional sensor data can contain redundant information that potentially hinders the efficiency and accuracy of clustering results. To address this issue, \textbf{Principal Component Analysis (PCA)} was applied to the data, reducing the feature space from 561 features to 50 principal components. 

The choice of 50 stems from the fact that, as there are over 500 features in the original data, 50 dimensions would be about 10\% of the original number of dimensions, significantly reducing computational complexity for the chosen clustering algorithms. Further, 50 components explain 87.1\% of the total variance in our data. As 80\%-90\% of total variance explained is generally considered an ideal range for principal components, the choice of reducing our data to 50 components strikes a balance between retaining a high percentage of total variance while reducing complexity.

As PCA is a linear dimensionality reduction technique, it was utilized as a baseline in our project. To test whether non-linear feature extraction would improve performance,  two neural network architectures were implemented to compress the 561-feature input into a 50-dimensional latent space, matching the PCA baseline for direct comparison. 

\subsubsection{Autoencoder (AE)} A symmetric autoencoder was implemented using TensorFlow and Keras [9]. The Encoder was constructed using a 150-dimensional dense hidden layer with ReLU activation and an additional dense layer that compresses the data into the 50-dimensional latent space. The Decoder mirrors the Encoder architecture in order to reconstruct the original input. The model was trained for 50 epochs using the Adam optimizer and Mean Squared Error (MSE) loss.

\subsubsection{Convolutional Variational Autoencoder (CVAE)} Also using TensorFlow and Keras [10][11], a CVAE was implemented to attempt to capture any temporal or spatial correlations in the data. The Encoder utilizes 1D Convolutional layers (filters are set to 32 and 64) followed by 1D MaxPooling layers to downsample the feature vector. The Encoder learns a probability distribution that a sampling layer utilizes to generate latent vectors. The Decoder mirrors the Encoder architecture in order to reconstruct the original input. The overall model attempts to minimize a total loss function that consists of Reconstruction Loss (MSE) and Kullback-Leibler (KL) Divergence. A small $\beta$ coefficient (0.0001) was applied to the KL term to reduce its weight on total loss, which allows for the balance of both reconstruction accuracy and latent space structure.

All of the clustering methods implemented in this project were tested on the PCA-processed data. After comparing the results, the highest performing clustering model, along with K-Means, was then tested on the two neural network architectures and compared.

\subsection{Clustering Methods}
As the primary goal of this project is a comparative analysis of standard unsupervised clustering approaches applied to a single dataset, 5 distinct clustering approaches were chosen based on the highest-performing algorithms found in the literature reviewed. For each of these methods, if a random state was able to be set, it was set to 42. Similarly, if a method required the number of clusters to be specified beforehand, 6 was used, as it is the number of distinct labels in our dataset.

\subsubsection{K-Means Clustering} This centroid-based approach was implemented first as a baseline [1][7]. Scikit-learn's \textbf{KMeans} was used with random initialization and $k=6$ clusters.

\subsubsection{Hierarchical Agglomerative Clustering} In order to investigate whether a connectivity-based approach would perform well on this data, we used Scikit-learn's \textbf{AgglomerativeClustering} [3]. Hierarchical clustering results were compared between four different linkage criteria: Single, Complete, Average, and Ward, all with the number of clusters specified to be 6 beforehand. 

\subsubsection{Spectral Clustering} To account for the potential of non-convex clusters, Scikit-learn's \textbf{SpectralClustering} was applied to the data [4]. Utilizing a nearest-neighbors affinity while specifying the number of clusters to be 6, clustering results using `kmeans' and `discretize' label assignments were compared.

\subsubsection{Density-Based Spatial Clustering of Applications with Noise (DBSCAN)} Following the work of some of the literature [2], Scikit-learn's \textbf{DBSCAN} was implemented to see if an approach that learns the number of clusters without any specification beforehand would perform well on the data. Parameters were tuned to achieve the best results, with $\text{eps}=11$ and $\text{min\_samples}=4$. 

\subsubsection{Gaussian Mixture Models (GMM)} Finally, in order to evaluate whether a probabilistic clustering approach would improve clustering performance [2][5], we applied Scikit-learn's \textbf{GaussianMixture}. Specifying the number of clusters to be 6, four covariance types were compared: Spherical, Diagonal, Tied, and Full.

\subsection{Evaluation Metrics}
In line with the four-phase unsupervised process described in the literature [1], these five methods were evaluated using both internal (data-only) and external (ground truth comparison) validation metrics.

\subsubsection{Silhouette Score} This is an internal metric that measures how well-separated the resulting learned clusters are. It assesses how similar a data point is to the cluster it is assigned to, as opposed to how similar it is to the other clusters, all without looking at the ground-truth labels. 

\subsubsection{Adjusted Rand Index (ARI)} As the ground-truth labels were provided in the UCI HAR dataset, ARI was utilized to measure how similar the clustering results were to the actual activity groupings. This external metric accounts for chance, meaning it is able to provide a more accurate measure of how well an unsupervised clustering model captures the true clusters.

After these two metrics were calculated for each variation of the chosen clustering approaches, the results were visualized via confusion matrices and scatter plots. The confusion matrices help to identify where any clustering overlap occurs between activity types. The 2D scatter plots use the first two principal components to aid in visualization and serve to assess cluster shapes and separability.

\section{Results and Discussion}

The methods for this project were executed in three major steps. The first involved \textbf{Data Processing, PCA, and K-Means}. The next step centered around \textbf{Clustering Method Comparisons} on the PCA-processed data. The final step was where the neural network architectures were implemented for the \textbf{Comparison of Dimensionality Reduction Techniques}.

\subsection{Data Processing, PCA, and K-Means}

This first step mainly served to assess whether or not a dimensionality technique like PCA would make a difference in unsupervised clustering performance. After processing the data, it was duplicated so that one of the duplicates could also be processed by PCA. Both the original copy of the data and the PCA-processed data were then given to a K-Means model. 

\begin{figure*}[!t]
\centerline{\includegraphics[width=\textwidth]{kmeans confusion matrices.png}}
\caption{K-Means Results on Original Data (left) vs. PCA Data (right). On both of the plots, the x-axis contains the ground truth activity labels, while the y-axis contains the learned K-Means cluster IDs. Each of the cells contains a value representing the number of data points from a given ground truth activity group that were predicted to belong to a particular learned cluster.}
\label{fig2}
\end{figure*}

The results of both clustering implementations are shown in \textbf{Fig. 2} in the form of confusion matrix heatmaps. It can be seen in Fig. 2 that, regardless of preprocessing type, K-Means does well in distinguishing between the (LAYING, SITTING, STANDING) and (WAKING, WALKING\_DOWNSTAIRS, WALKING\_UPSTAIRS) activity groups. However, it has trouble distinguishing activity types within either of the groups. 

It can also be seen in confusion matrices that there isn't much difference between the clustering results of K-Means on the original data vs. the results of K-Means on the PCA-processed data. This lack of difference can also be observed in the resulting Adjusted Rand Index (ARI) and Silhouette scores. K-Means on the original data achieved an ARI Score of \textbf{0.4196} and a Silhouette Score of \textbf{0.1099}, while K-Means on the PCA-processed data achieved an ARI Score of \textbf{0.4201} and a Silhouette Score of \textbf{0.1524}. 

The scores from the implementation that included PCA are slightly higher than the ones achieved without, all while having a significantly smaller number of features perform calculations on during the clustering process (50 vs. 561). Because of this,  PCA-processed data was utilized with all of the clustering implementations present in the following subsection.

\subsection{Clustering Method Comparisons}

This step involved the comparison of five traditional unsupervised clustering algorithms applied to the 50-dimensional PCA baseline. In addition to \textbf{K-Means}, the algorithms implemented in this step are \textbf{Hierarchical Agglomerative Clustering}, \textbf{Spectral Clustering}, \textbf{DBCAN}, and \textbf{GMM}. When determining whether a given clustering algorithm outperformed another, the ARI score was prioritized due to it being a measure of how similar the learned clusters are to the ground truth.

As the scores from the K-Means implementation on the PCA-processed data were reported in the previous subsection, we will only briefly go over the results of the other four algorithms in the following list:

\subsubsection{Hierarchical Clustering} Out of the four linkage types implemented, \textbf{Ward} performed the best (ARI = \textbf{0.4218}), followed by Complete (ARI = 0.0794), then Average (ARI = 0.0032), and finally Single (ARI = 0.0001).

\subsubsection{Spectral Clustering} Between the two label assignment types implemented, \textbf{`discretize'} achieved a slightly higher ARI score of \textbf{0.4947} compared to the `kmeans' ARI score of 0.4922.

\subsubsection{DBSCAN} Using the tuned parameters of $\text{eps}=11$ and $\text{min\_samples}=4$, DBSCAN estimated a number of 34 clusters, estimated 1,791 points to be noise out of 10,299, and achieved an ARI score of \textbf{0.3223}.

\subsubsection{GMM} Comparing the four covariance types implemented, \textbf{Tied} achieved the highest ARI score (ARI = \textbf{0.4278}), followed by Full (ARI = \textbf{0.3181}), then Diagonal (ARI = \textbf{0.3015}), and finally Spherical (ARI = \textbf{0.2765}).

\begin{figure}[!t]
\centerline{\includegraphics[width=\columnwidth]{clustering scores.png}}
\caption{Comparison Table of Clustering Model Performances on the PCA-Processed Data. The rows are sorted by ARI score in descending order. Spectral achieved the highest ARI score (highlighted in green) while DBSCAN resulted in the lowest ARI score (highlighted in magenta).}
\label{fig3}
\end{figure}

A table summarizing the scores from the best implementations of each of the five clustering algorithms is shown in \textbf{Fig. 3}. As can be seen in the table, Spectral Clustering (with `discretize' labels) achieved the highest Adjusted Rand Index (ARI). This suggests that the relationship between activity groups in the dataset is likely non-convex, meaning an affinity-based approach like Spectral Clustering would outperform an approach that assumes spherical clusters, such as K-Means.

Among the variations of Hierarchical implementations, the choice of linkage proved to greatly impact clustering results, as Ward linkage outperformed the others by a significant amount. Additionally, the Tied covariance GMM outperformed the covariance types, meaning the clusters are likely to share a similar orientation to each other in the feature space.

DBSCAN demonstrated poor suitability for this dataset, as it estimated 34 clusters and labeled approximately 17\% of the data as noise even after fine-tuning the parameters to maximize performance scores. This can be an indication that the true activity clusters vary significantly in density, a challenge for the DBSCAN approach implemented in this project.

\begin{figure*}[!t]
\centerline{\includegraphics[width=\textwidth]{comparison scatter plot.png}}
\caption{Clustering Visualization for Model Comparison. Each of the scatter plots utilizes the first two principal components for 2D visualization, with the first principal component (PC1) on the x-axis and the second principal component (PC2) on the y-axis. The `Actual Activity Labels' plot contains a legend of ground-truth cluster labels for reference. The rest of the plots don't contain a legend as the learned clusters are represented by IDs (i.e., 0-5).}
\label{fig4}
\end{figure*}

\textbf{Fig. 4} provides a visualization of clustering results from each of the models, as well as the ground truth clusters. It isn't difficult to see how the clusters in the Spectral scatter plot appear more visually similar to the Actual Activity clusters than the clusters in the rest of the scatter plots.

\subsection{Comparison of Dimensionality Reduction Techniques}

This final step is where we implemented a basic Autoencoder (AE) and a Convolutional Variational Autoencoder (CVAE) in order to obtain two different 50-dimensional learned latent spaces. Clustering was then performed on the data taken out of this latent space. As Spectral (Discretize) proved to perform the best on the PCA-processed data, it was tested on the latent space data along with K-Means. 

\begin{figure}[!t]
\centerline{\includegraphics[width=\columnwidth]{dr scores.png}}
\caption{Comparison Table of Clustering Performance on Various Dimensionality Reduction Techniques. Clustering was performed on PCA-processed data, data from the Autoencoder (AE) latent space, and data from the Convolutional Variational Autoencoder (CVAE) latent space. K-Means on the CVAE-processed data achieved the highest ARI score (highlighted in green) while K-Means on the PCA-processed data resulted in the lowest ARI score (highlighted in magenta). Note that due to the nature of our CVAE implementation, we can only achieve approximate values for the scores even after setting the random seed to 42.}
\label{fig5}
\end{figure}

In \textbf{Fig. 5}, we can see that overall, the integration of neural network-based dimensionality reduction improved clustering results over the linear PCA. The \textbf{AE + Spectral} implementation achieved an ARI score of \textbf{0.519}, an approximately 4.8\% improvement over the PCA + Spectral version.

However, what is more surprising and significant are the results of the CVAE implementations. While the CVAE + Spectral implementation performed worse than AE + Spectral with an ARI score of 0.496, the \textbf{CVAE + K-Means} clustering implementation achieved an ARI score of approximately \textbf{0.650}, outperforming all other dimensionality reduction/clustering model combinations implemented in this project. This improvement is an approximately 54\% increase in clustering accuracy compared to the PCA + K-Means baseline.

\begin{figure*}[!t]
\centerline{\includegraphics[width=\textwidth]{autoencoder scatter plots.png}}
\caption{Clustering Visualization for Model Comparison in AE Latent Space. Each of the scatter plots utilizes the first two principal components of the latent space for 2D visualization, with the first principal component (PC1) on the x-axis and the second principal component (PC2) on the y-axis. The `Ground Truth' plot contains a legend of ground-truth cluster labels for reference.}
\label{fig6}
\end{figure*}

\begin{figure*}[!t]
\centerline{\includegraphics[width=\textwidth]{cvae scatter plot.png}}
\caption{Clustering Visualization for Model Comparison in CVAE Latent Space. Each of the scatter plots utilizes the first two principal components of the latent space for 2D visualization, with the first principal component (PC1) on the x-axis and the second principal component (PC2) on the y-axis. The `Ground Truth' plot contains a legend of ground-truth cluster labels for reference.}
\label{fig7}
\end{figure*}

Similar to Fig. 4, \textbf{Fig. 6} and \textbf{Fig. 7} serve to visualize clustering via scatter plots. Fig. 6 compares K-Means and Spectral clustering against the ground truth in the AE latent space, while Fig. 7 compares them in the CVAE latent space. It is easy to see how the Spectral clusters are more visually similar to the actual clusters in the AE plots, while the K-Means clusters are more similar to the actual clusters in the CVAE plots.

\subsection{Discussion of Metrics}

A detail worth noting in the results of this project is the seemingly inverse relationship between the ARI and Silhouette scores. This can be seen in particular when looking at the CVAE rows of \textbf{Fig. 5}. There we see that the CVAE + K-Means model produced a relatively low Silhouette Score (approximately 0.062) despite achieving the highest ARI score (0.650). 

This inverse relationship in scores suggests that while the CVAE was able to successfully group similar activities in the latent space (as supported by the high ARI score), the learned clusters are likely overlapping, which leads to a lower internal separation (Silhouette) score. This can be seen visually in the middle scatter plot of \textbf{Fig. 7}.

The reason for the superior performance of the CVAE over the AE is likely the inclusion of the 1D Convolutional layers. These layers allow the CVAE to better capture any spatial-temporal patterns present in the features of the original sensor data. This, in turn, allows for the creation of a more meaningful latent representation for physical movement compared to any latent representation that is created by a standard linear or dense layer present in a basic AE.

\section{Future Work and Conclusions}

\subsection{Future Work}
Building on the results of this project, there are several possible next steps. 

\begin{itemize}
    \item \textbf{Further Model Testing on AE/CVAE:} As the utilization of an Autoencoder (AE) and Convolutional Variational Autoencoder (CVAE) was added on as extensions in this project, only K-Means and Spectral Clustering were tested on the data processed by the neural networks. As it was unexpected to us that K-Means would end up outperforming Spectral when clustering data from the CVAE latent space, it would be interesting to know if the other models tested on the PCA-processed data would also perform better on the CVAE-processed data.
    \item \textbf{Hyperparameter Tuning:} The AE and CVAE architectures in this project utilized baseline hidden layer dimensions and learning rates based on example implementations [9][10][11]. Further, the dimension of the latent space was chosen to match the previous selection of 50 PCA features. Utilizing known hyperparameter tuning methods to obtain the optimal values for the layer dimensions, the $\beta$ coefficient (for KL loss), and the latent dimension could further improve clustering performance on HAR data.
    \item \textbf{Additional Feature Processing:} Utilizing techniques discussed in the literature, such as Correlation Feature Selection (CFS) [3] or focusing on Acceleration Vector Magnitude (AVM) [5], before passing data into neural network encoders may help to reduce noise and improve the stability of the latent representations.
    \item \textbf{Self-Supervised Hybrid Models:} Exploring more state-of-the-art clustering architectures like Semantic Clustering by Adopting Nearest Neighbors (SCAN) [6] could also improve HAR data clustering results. The CVAE + K-Means approach that we found to perform the best is a more primitive example of one of these hybrid frameworks. Thus, more advanced hybrid models have the potential to push ARI scores even higher on this data.
\end{itemize}

\subsection{Conclusion}

This project compared a diverse range of unsupervised clustering methods and dimensionality reduction techniques using the UCI Human Activity Recognition dataset. The results demonstrate that while traditional clustering algorithms like Spectral Clustering and Gaussian Mixture Models can provide a respectable baseline on data processed by the linear Principal Component Analysis (PCA), these approaches are ultimately limited by their inability to capture the non-linear, spatial-temporal relationships present in data gathered from wearable sensors.

The primary takeaway of this project is the significant performance gain offered when using neural networks. The transition from PCA data features to Convolutional Variational Autoencoder (CVAE) latent space features resulted in an Adjusted Rand Index (ARI) score increase from 0.420 to approximately 0.650. This suggests that the 1D Convolutional layers in our CVAE implementation were able to better capture the non-linear relationships and patterns that are present in physical activity data.

Additionally, this project highlights the discrepancy between internal metrics, such as the Silhouette scores, and external metrics, such as the ARI scores. The finding that the CVAE-based implementations achieved higher activity grouping accuracy despite lower cluster separation indicates that latent spaces for Human Activity Recognition (HAR) data learned by state-of-the-art models may consist of dense and overlapping clusters, which distance-based metrics may fail to characterize correctly.

In conclusion, while deep learning approaches are at the forefront of modern HAR research, the choice of a robust dimensionality reduction framework can keep traditional unsupervised clustering methods relevant, allowing for the achievement of high accuracy in real-world human activity recognition.

\begin{thebibliography}{00}
\normalsize

\bibitem{b1} P. Ariza Colpas, E. Vicario, E. De-La-Hoz-Franco, M. Pineres-Melo, A. Oviedo-Carrascal and F. Patara, ``Unsupervised Human Activity Recognition Using the Clustering Approach: A Review,'' \emph{Sensors,} vol. 20, no. 9, 2020, doi:10.3390/s20092702.

\bibitem{b2} Y. Kwon, K. Kang and C. Bae, ``Unsupervised learning for human activity recognition using smartphone sensors,'' \emph{Expert Systems with Applications,} vol. 41, no. 14, 2014, doi:10.1016/j.eswa.2014.04.037.

\bibitem{b3} C. Dobbins and R. Rawassizadeh, ``Towards Clustering of Mobile and Smartwatch Accelerometer Data for Physical Activity Recognition,'' \emph{Informatics,} vol. 5, no. 2, 2018, doi:10.3390/informatics5020029.

\bibitem{b4} V. Vardhan Baligodugula and F. Amsaad, ``Unsupervised Learning: Comparative Analysis of Clustering Techniques on High-Dimensional Data,'' 2025, doi:10.48550/arXiv.2503.23215.

\bibitem{b5} K. Jurcic and R. Magjarevic, ``Time Domain Cluster Analysis of Human Activity Using Triaxial Accelerometer Data,'' \emph{CMBES Proc.,} vol. 46, Jun. 2024.

\bibitem{b6} A. Ahmed, H. Haresamudram and T. Ploetz, ``Clustering of Human Activities from Wearables by Adopting Nearest Neighbors,'' \emph{Proceedings of the ACM ISWC,} New York, NY, USA, 2022, doi:10.1145/3544794.3558477.

\bibitem{b7} M. Sitorus, M. Melisa and D. Herninda, ``Driver Predicting Behavior Based on Accelerometer and Gyroscope Sensors with K-Means Algorithm Method,'' \emph{RIGGS,} vol. 2, no. 1, pp. 12-17 May 2023, doi:10.31004/riggs.v2i1.29.

\bibitem{b8} J. Reyes-Ortiz, D. Anguita, A. Ghio, L. Oneto and X. Parra, ``Human Activity Recognition Using Smartphones,'' \emph{UCI Machine Learning Repository,} 2013. doi:10.24432/C54S4K.

\bibitem{b9} mrgrhn. ``AutoEncoders with TensorFlow.'' Medium. https://medium.com/analytics-vidhya/autoencoders-with-tensorflow-2f0a7315d161 (accessed April 13, 2026).

\bibitem{b10} S. Sekhar Bhakta. ``Convolutional Variational Autoencoder in Tensorflow.'' GeeksforGeeks. https://www.geeksforgeeks.org/deep-learning/convolutional-variational-autoencoder-in-tensorflow/ (accessed April 15, 2026).

\bibitem{b11} TensorFlow Staff. ``Convolutional Variational Autoencoder.'' TensorFlow. https://www.tensorflow.org/tutorials/generative/cvae (accessed April 15, 2026).

\end{thebibliography}

\end{document}
